\textbf{16.} Докажите $\textbf{PSPACE}$-полноту языка $\mathsf{LOOP} = \{(M ,x, 1^s)\colon \text{вычисление $M(x)$ зацикливается,}$\\$\text{заняв не более $s$ ячеек на ленте}\}$.

\textbf{Решение.} Как это обычно бывает с полнотой, сначала докажем принадлежность целевому классу, а потом --- трудность задачи.

1. $\mathsf{LOOP} \in \mathbf{PSPACE}$.\\
Пусть на вход даётся тройка $(M, x, 1^s)$. Машина Тьюринга, распознающая язык $\mathsf{LOOP}$, запустит машину $M$ на входе $x$ на $a^{s}q|x|s$ шагов (где $a$ -- размер алфавита, $q$ -- число состояний $M$) и будет внимательно следить, сколько памяти она кушает. Число шагов мы взяли именно таким, чтобы оно было больше возможного числа конфигураций $M$. Если машина $M$ попытается использовать больше $s$ ячеек памяти, то мы сразу останавливаемся и переходим в $q_{reject}$. Аналогичный исход нас ожидает, если $M$ остановилась --- мы же хотим, чтобы она зациклилась. А вот если она не завершила вычисление после установленного числа шагов, то мы можем гарантировать, что она уже никогда не остановится: по принципу Дирихле какая-то конфигурация встречалась дважды, а значит она будет повторяться вновь и вновь.

Такая машина будет тратить лишь полиномиальное от входа число ячеек: ей нужно порядка $s$ ячеек памяти для симулирования машины $M$, счётчик для хранения числа используемых $M$ ячеек памяти (оно не больше $s$, иначе стоп машина) и счётчик для хранения номера текущего шага (шагов не больше $2^{poly(s)}$, а запись этого числа требует $poly(s)$ памяти) -- всё это является полиномом от длины входа.

2. $\mathsf{LOOP}$ --- $\mathbf{PSPACE}$-трудный.\\
Пусть $A$ --- произвольный язык из $\mathbf{PSPACE}$. По определению этого класса для $A$ существует такая машина Тьюринга $M$, что вычисление $M(x)$ требует не более, чем $c\cdot n^d$ ячеек памяти, где $n = |x|$, $c$ и $d$ --- какие-то константы. По этой машине построим новую машину $M'$: по входу $x$ она моделирует вычисление $M(x)$, и если она вернула 1, то мы входим в цикл, иначе -- возвращаем что-нибудь. Хочется доказать, что $x \in A \Longleftrightarrow (M', x, 1^{c\cdot n^d}) \in \mathsf{LOOP}$.

Машине $M'$ всегда хватит $c\cdot n^d$ ячеек памяти, так как она просто моделирует машину $M$, которой этого количества достаточно, поэтому принадлежность $(M', x, 1^{c\cdot n^d})$ к $\mathsf{LOOP}$ равносильна зацикливанию $M'(x)$, что по построению равносильно $x \in A$. Также несложно доказывается, что такая сводимость полиномиальная: код для машины $M'$ постоянный и не зависит от $x$, само число $x$ переписать можно очень быстро, а написать $c\cdot n^d$ единиц требует полиномиального времени.